\section{The mechanism that could transfer}

The proof is not merely an analytic manipulation of Dirichlet
$L$-functions. Its main comparison is between upper and lower bounds for
the same arithmetic size. The public guide and the exact Lean statements
identify three inputs \cite{guide,mainineq}:

\begin{enumerate}
\item \textbf{Short forms.} Theta derivatives have local divisibility that
records smoothed prime sums. After a controlled normalization, suitable
integral source forms map to short target vectors with bounded index.
\item \textbf{Enough forms.} Near-isometry and a second-moment estimate
show that few translates of the source map span almost all of the target.
This is where the quantitative improvement below applies.
\item \textbf{Long integral forms.} Determinants of chains of lifts,
regular families of twists, and bounds on denominators provide a lower
bound on the size of sufficiently many integral target forms.
\end{enumerate}

The common leading contribution cancels when these bounds are combined.
The conclusion of \src{New/Reduction/MainInequality.lean}
{\texttt{New.star\_inequality}} is
\[
 m\,\mathrm{totalRemainder}\le f\,a(\log b)^5.
\]
Here $m$ is the number of successive lifts, $f$ is the number of infinite
places, $a$ is a derivative order, and $b$ is a much larger auxiliary
weight. The source chooses $\sqrt{x}\le m\le2\sqrt{x}$ and several widely
separated scales. The main inequality and a prime number theorem at a
second cutoff imply a one-sided estimate for
\[
 R_{d,q}(x)=
 \sum_{\substack{n\le x\\n\equiv1\pmod q}}
 \Lambda(n)(1-\chi_d(n))\left(\frac1n-\frac1x\right)
 -\frac{\log x}{\varphi(q)}-c_{d,q}.
\]
For the specified auxiliary primes $d$ and moduli $q$, this estimate is
$R_{d,q}(x)\le x^{-1/2+\eps}$ eventually, for every $\eps>0$
\cite{onesided}.

\subsection{Why the last analytic step is not the main missing extension}

The transform of this remainder has residues given by differences of
zero multiplicities of Dirichlet functions and their quadratic twists.
The source chooses an auxiliary twist that is zero-free in a prescribed
rectangle. This leaves the relevant residues with the sign needed for
Landau's one-sided argument. The declaration
\src{New/Reduction/Zeros/Criterion.lean}
{\texttt{New.LFunction\_ne\_zero\_of\_primeRemainder\_le}}
makes the quantifiers over the twists and the eventual prime-sum bound
explicit \cite{criterion}.

For another family, a finite-group identity can identify which prime
statistic ought to detect its zeros. It does not prove the corresponding
prime-sum estimate. Nor does it automatically eliminate real poles or
cancellation by auxiliary twists. This distinction is central to the
research assessment: numerous previous reports have already constructed
the detectors, while the required integral realization remains missing.

\subsection{Where a quantitative improvement has leverage}

In a prospective extension, a smaller target span causes two costs.
It can enlarge error terms after division by the selected rank, and it
can increase the auxiliary completion needed to compare partial and full
determinants. Reducing both the missing rank and the number of translates
therefore gives more room for new denominator and ramification losses.

This is a benefit to a future construction, not a proof that such a
construction exists. In particular, improving $b^{-2}$ to $b^{-4}$ in the
rank loss does not replace the square-root scale $m\asymp\sqrt{x}$, prove
an exponent beyond $1/2$ for the prime error, or control the spacing or
multiplicity of zeros on the critical line.

There is also a simple reason to keep the rank information explicit.
Two positive forms can have the same full determinant and arbitrarily
different determinants on a fixed line: the diagonal matrices
$\operatorname{diag}(M,M^{-1})$ and $I_2$ have determinant ratio one,
but ratio $M$ on the first coordinate line. A full determinant identity
alone therefore cannot repair an uncontrolled missing subspace.
